\chapter{Flow Market Making at a Bank}\label{ch:s2:flow-market-making-at-a-bank}

A bank's flow desk makes money by pricing clients, not by predicting prices. It learns which clients' trades will move against it and charges them
accordingly, and it keeps the risk of one client's trade until another client's offsets it. Butz and Oomen estimated that a representative
large dealer takes several minutes on average to internalise a client's trade in the most liquid currencies, and tens of minutes in emerging markets.
On this chapter's synthetic desk, pricing each client from its own history raises the \omterm{def:s2:flow-market-making-at-a-bank:franchise}{franchise P\&L} from \$6.63 million at the best
of four flat prices to \$7.60 million. Uninformed clients pay 1.39 basis points and informed ones stop trading. A small axe skew raises the share of risk
internalised from 91\% to 97\%. The build is \texttt{firm.flowmm}.

\section{The franchise and its P\&L}

\begin{definition}[Franchise P\&L]\label{def:s2:flow-market-making-at-a-bank:franchise}
The \emph{franchise P\&L}\index{franchise P\&L} of a flow desk is its profit from serving clients, decomposed into spread capture (the price clients
pay relative to mid), the marks of the inventory their trades leave (where adverse selection shows), and the cost of hedging what it does not
internalise.
\end{definition}

\texttt{firm.flowmm} sends a desk 200\,000 requests for quote from 90 clients, each for one unit of \$1 million, so that a basis point is \$100. Sixty
clients are uninformed: after they trade, the mid moves at random. Twenty are moderately informed: the mid drifts 0.8 basis points in their direction
over the next ten requests. Ten are informed, with a drift of 2.5 basis points. The drift happens whoever fills the trade, so the mid path is the same
whatever the desk quotes, and every pricing policy is compared on the same prices. Each client also has its own competing quote, the half-spread it
can get elsewhere: 1.0 to 2.5 basis points for uninformed clients, some of them captive, 0.8 to 1.5 for the moderate ones and 0.4 to 0.8 for the
informed ones, who shop around. A client takes the desk's price with a probability that falls smoothly as the price rises above its competing quote
(\cref{lst:s2:flow-market-making-at-a-bank:desk}).

The desk quotes everyone 1.2 basis points for the first 100\,000 requests and learns from them. The mark-out, the mid's move over the next 20
requests against the desk's side, separates the three types (\cref{fig:s2:flow-market-making-at-a-bank:markout}): $-0.01$, 0.84 and 2.41 basis points
on average. The hit ratio, the share of a client's requests the desk wins, reveals its competing quote: 80\%, 43\% and 8\%. Inferred from each
client's hit ratio, the competing quotes have a correlation of 0.998 with the true ones.

\begin{omfigure}
\begin{tikzpicture}
\begin{axis}[width=11cm,height=5.2cm,xlabel={\small requests after the trade},ylabel={\small mark-out (bp, client's favour)},
  tick label style={font=\footnotesize},xmin=0,xmax=40,ymin=-0.5,ymax=3,grid=major,grid style={gray!25},table/col sep=comma,
  legend style={font=\scriptsize,at={(0.5,-0.3)},anchor=north},legend columns=3]
\addplot[black,thick] table[x=horizon,y=uninformed]{figdata/strategies-2/24-flow-market-making-at-a-bank/markout.csv};
\addplot[omDef,thick,dashed] table[x=horizon,y=moderate]{figdata/strategies-2/24-flow-market-making-at-a-bank/markout.csv};
\addplot[omThm,thick] table[x=horizon,y=informed]{figdata/strategies-2/24-flow-market-making-at-a-bank/markout.csv};
\legend{uninformed,moderately informed,informed}
\end{axis}
\end{tikzpicture}

\omcaption{Mark-out curves of the synthetic desk's filled trades by client type: the mid's move after the trade in the client's favour, in basis
points, against the number of requests since. Data: \texttt{s2\_flowmm.markout\_curve}.}\label{fig:s2:flow-market-making-at-a-bank:markout}
\end{omfigure}

\section{Client tiering}

\begin{definition}[Client tiering]\label{def:s2:flow-market-making-at-a-bank:tiering}
\emph{Client tiering}\index{client tiering} is pricing each client, or group of clients, by what its trades cost the desk (its mark-out) and by how
sensitive it is to price (its hit ratio against the desk's quotes), rather than quoting everyone the same.
\end{definition}

For each client the desk picks the half-spread that maximises its expected profit per request, the half-spread less the client's mark-out, times the
chance of winning the request at that price (\cref{lst:s2:flow-market-making-at-a-bank:price}). Captive uninformed clients pay more, price-sensitive
ones less, and informed clients are quoted out of the market. On the second 100\,000 requests (P\&L in \$ millions):

\begin{center}
\small
\begin{tabular}{@{}lccccc@{}}
\toprule
policy & trades & capture & marks & hedging & \omterm{def:s2:flow-market-making-at-a-bank:franchise}{franchise P\&L}\\
\midrule
flat 1.0 bp & 74\,622 & 7.46 & $-1.21$ & 0.35 & 5.90\\
flat 1.2 bp & 64\,307 & 7.72 & $-0.84$ & 0.31 & 6.57\\
flat 1.4 bp & 53\,310 & 7.46 & $-0.58$ & 0.25 & 6.63\\
flat 1.6 bp & 42\,399 & 6.78 & $-0.36$ & 0.20 & 6.23\\
priced by client & 62\,147 & 8.53 & $-0.65$ & 0.28 & 7.60\\
priced by client, axe skew 0.05 & 61\,519 & 8.35 & $-0.62$ & 0.10 & 7.63\\
\bottomrule
\end{tabular}
\end{center}

No flat price does as well. A low one wins the informed clients' trades and loses on them; a high one loses the captive clients' business along with
the informed. By type, from the flat 1.2 basis points to prices by client:

\begin{center}
\small
\begin{tabular}{@{}lcccc@{}}
\toprule
per trade (bp) & trades & capture & mark-out & net\\
\midrule
uninformed, flat & 53\,500 & 1.20 & 0.00 & 1.20\\
uninformed, priced & 53\,600 & 1.39 & 0.00 & 1.40\\
moderate, flat & 9\,869 & 1.20 & 0.83 & 0.37\\
moderate, priced & 8\,537 & 1.25 & 0.82 & 0.44\\
informed, flat & 938 & 1.20 & 2.38 & $-1.18$\\
informed, priced & 10 & & &\\
\bottomrule
\end{tabular}
\end{center}

Oomen showed that price signatures, the average price path around executions, can tell internalising liquidity providers from externalising ones;
the same curves, computed by client, are what a desk tiers on. Tiering has a mirror image: a client's execution quality depends on how its dealers
manage risk. Oomen's study of last look, the dealer's option to reject a deal request after seeing where the price has gone, found that the design
of that option determines the client's execution risk but need not change its effective costs.

\section{Axes and skewing}

\begin{definition}[Axe-driven skew]\label{def:s2:flow-market-making-at-a-bank:skew}
An \emph{axe-driven skew}\index{axe-driven skew} is a shift of both sides of a dealer's quote against its inventory or its axes, so that clients are
offered better prices on the trades that reduce the desk's position and worse ones on those that add to it.
\end{definition}

The desk skews its quotes by 0.05 basis points per unit of inventory: long five units, it sells 0.25 basis points cheaper and buys 0.25 basis
points cheaper too. Clients then trade more in the direction that reduces the position, and the desk hedges less. The skew costs capture and saves
hedging and risk (\cref{fig:s2:flow-market-making-at-a-bank:skew}): at 0.05 the \omterm{def:s2:flow-market-making-at-a-bank:franchise}{franchise P\&L} is slightly higher, \$7.63 million, and the average
inventory falls from 2.75 to 2.18 units; at 0.2 the desk never hedges, holds 1.30 units on average and earns \$7.08 million.

\begin{omfigure}
\begin{tikzpicture}
\begin{axis}[width=11cm,height=5.2cm,axis y line*=right,axis x line=none,ylabel={\small share internalised (\%)},
  tick label style={font=\footnotesize},xmin=0,xmax=0.2,ymin=88,ymax=101,table/col sep=comma]
\addplot[omDef,thick,dashed,mark=square*,mark size=1pt] table[x=skew,y=internalised]{figdata/strategies-2/24-flow-market-making-at-a-bank/skew.csv};
\label{plot:s2:flowint}
\end{axis}
\begin{axis}[width=11cm,height=5.2cm,axis y line*=left,xlabel={\small axe skew (bp per unit of inventory)},
  ylabel={\small franchise P\&L (bp thousands)},tick label style={font=\footnotesize},xmin=0,xmax=0.2,ymin=70,ymax=77,grid=major,
  grid style={gray!25},table/col sep=comma,scaled x ticks=false,xticklabel style={/pgf/number format/fixed},
  legend style={font=\scriptsize,at={(0.5,-0.3)},anchor=north},legend columns=2]
\addplot[omThm,thick,mark=*,mark size=1pt] table[x=skew,y=net]{figdata/strategies-2/24-flow-market-making-at-a-bank/skew.csv};
\addlegendentry{franchise P\&L (left)}
\addlegendimage{/pgfplots/refstyle=plot:s2:flowint}\addlegendentry{share internalised (right)}
\end{axis}
\end{tikzpicture}

\omcaption{The synthetic desk priced by client, second half: \omterm{def:s2:flow-market-making-at-a-bank:franchise}{franchise P\&L} (in thousands of basis points of unit notional, \$100 each) and the share
of client volume internalised, against the axe skew. Data: \texttt{s2\_flowmm.policy}.}\label{fig:s2:flow-market-making-at-a-bank:skew}
\end{omfigure}

\section{Internalisation}

A desk internalises a client's trade when it keeps the risk until another client's trade offsets it, and externalises it when it hedges in the
market. The synthetic desk holds up to five units and hedges beyond them at the market's half-spread of 0.5 basis points. Priced by client without a
skew, it internalises 90.8\% of its volume; with the 0.05 skew, 96.8\%. Butz and Oomen modelled internalisation as a queue and found that its costs
are lower for dealers willing to hold more risk and for those whose clients are more price-sensitive. They concluded that a client's costs depend on
the dealer's risk management as well as its own trading, and that clients should tell externalisers from passive and aggressive internalisers.

\section{Strategy files}

\begin{strategyfile}{Tiered pricing by client mark-out}\label{strat:s2:flow-market-making-at-a-bank:tier}
\sfield{Who pays you, and why} Clients who value the relationship, convenience or credit over the last basis point.
\sfield{Instruments and venues} Requests for quote and streaming prices to clients, by platform or voice.
\sfield{Signal} Each client's mark-out curve and hit ratio.
\sfield{Sizing and execution} Half-spread by client that maximises expected profit per request.
\sfield{Costs} Adverse selection; lost business from over-pricing.
\sfield{How it dies} Clients who change behaviour; competitors who tier better.
\sfield{Horizon, capacity, infrastructure} Continuous; client analytics and pricing engines.
\sfield{Backtest honestly} Mark-outs on trades won only; the win curve estimated from varied quotes, not assumed.
\sfield{Sources} Oomen (2019); this chapter: \$7.60 million against \$6.63 million at the best flat price.
\end{strategyfile}

\begin{strategyfile}{Axe skewing}\label{strat:s2:flow-market-making-at-a-bank:skew}
\sfield{Who pays you, and why} Clients on the other side of the desk's position, who get a better price.
\sfield{Instruments and venues} The desk's quotes; axes shown to clients.
\sfield{Signal} Inventory and axes.
\sfield{Sizing and execution} A skew per unit of inventory, set against capture lost and hedging saved.
\sfield{Costs} Capture given up.
\sfield{How it dies} Clients who read the skew and trade against it.
\sfield{Horizon, capacity, infrastructure} Continuous.
\sfield{Backtest honestly} Clients' response to the skew estimated, not assumed.
\sfield{Sources} This chapter: internalisation from 90.8\% to 96.8\% at 0.05.
\end{strategyfile}

\begin{strategyfile}{Internalisation before hedging}\label{strat:s2:flow-market-making-at-a-bank:internalise}
\sfield{Who pays you, and why} The market's spread not paid on offsetting flows.
\sfield{Instruments and venues} The desk's own inventory; the market to hedge the residual.
\sfield{Signal} Expected offsetting flow against the risk of waiting.
\sfield{Sizing and execution} A risk limit beyond which the desk hedges.
\sfield{Costs} Inventory risk; information leaked by hedging.
\sfield{How it dies} One-way markets.
\sfield{Horizon, capacity, infrastructure} Minutes to hours; a central risk book.
\sfield{Backtest honestly} Inventory marks over the internalisation horizon.
\sfield{Sources} Butz and Oomen (2019).
\end{strategyfile}

\begin{strategyfile}{RFQ win-rate management}\label{strat:s2:flow-market-making-at-a-bank:winrate}
\sfield{Who pays you, and why} Clients whose business is worth winning at a price.
\sfield{Instruments and venues} Multi-dealer request-for-quote platforms.
\sfield{Signal} Hit ratios by client, size and time against quoted spreads.
\sfield{Sizing and execution} Target win rates by client segment.
\sfield{Costs} Winner's curse on informed requests.
\sfield{How it dies} Chasing volume at negative mark-out.
\sfield{Horizon, capacity, infrastructure} Continuous; platform analytics.
\sfield{Backtest honestly} Lost requests' outcomes are unobserved; do not assume them.
\sfield{Sources} No performance figure verified; this chapter's hit ratios: 80\%, 43\% and 8\%.
\end{strategyfile}

\section{Tutorial: pricing the client}\label{tut:s2:flow-market-making-at-a-bank:tut}

\begin{tutorial}
\textbf{Goal.} Simulate a desk facing clients with different information and competition, learn their mark-outs and hit ratios, price them, skew to
the inventory and decompose the \omterm{def:s2:flow-market-making-at-a-bank:franchise}{franchise P\&L}. \textbf{End state:} the three tables and three figures.
\begin{tutsteps}
\item \textbf{The desk}.
\omcode{code/firm/flowmm/firm_flowmm.py}{63}{94}{Quote, fill, internalise, hedge beyond the limit, decompose.}\label{lst:s2:flow-market-making-at-a-bank:desk}
\item \textbf{Learning and pricing the clients}.
\omcode{code/firm/flowmm/firm_flowmm.py}{97}{130}{Mark-outs, hit ratios and the price that maximises expected profit.}\label{lst:s2:flow-market-making-at-a-bank:price}
\item \textbf{Run} \texttt{estimates()}, \texttt{policy()}, \texttt{by\_type()} and \texttt{fig\_flowmm.py}.
\end{tutsteps}
\textbf{What to change next.} Let clients change type halfway; estimate the win curve from randomised quotes instead of assuming its shape; add last
look.
\end{tutorial}

\section{Build: flow market making}\label{bld:s2:flow-market-making-at-a-bank:flowmm}

\begin{build}
\bfield{Purpose} A request-for-quote desk with client mark-outs, pricing by client, axe skewing and internalisation, with its P\&L decomposed.
\bfield{Interface} \texttt{FlowConfig(\dots)}, \texttt{simulate\_flow(cfg)}, \texttt{run\_desk(flow, cfg, half\_spread, skew, start, end)},
\texttt{client\_markouts}, \texttt{client\_hits}, \texttt{price\_clients}.
\bfield{Rules} The mid path does not depend on the desk; the client's choice is logistic in the price paid; hedging only beyond the limit.
\bfield{Acceptance tests} \texttt{code/firm/flowmm/tests/}: the decomposition adds up and wider quotes fill less; mark-outs rank the types; pricing
by hand; inventory within the limit.
\bfield{Stretch} Last look; streaming prices; several products with correlated inventory.
\end{build}

\begin{omsources}
\item M.~Butz and R.~Oomen, ``Internalisation by electronic FX spot dealers'', \emph{Quantitative Finance} 19(1), 2019.
\item R.~Oomen, ``Last look'', \emph{Quantitative Finance} 17(7), 2017.
\item R.~Oomen, ``Price signatures'', \emph{Quantitative Finance} 19(5), 2019.
\end{omsources}

\section{Exercises}

\begin{exercise}[$\star$]\label{exo:s2:flow-market-making-at-a-bank:1}
A client pays 1.2 basis points on \$1 million and its trades mark out at 0.83. What does the desk earn per trade, in dollars?
\end{exercise}

\begin{exercise}[$\star$]\label{exo:s2:flow-market-making-at-a-bank:2}
The desk is long 4 units with a skew of 0.05 basis points per unit and a half-spread of 1.2. What does a buying client pay, and a selling one?
\end{exercise}

\begin{exercise}[$\star$]\label{exo:s2:flow-market-making-at-a-bank:3}
A client's hit ratio at 1.2 basis points is 80\% and the choice has a width of 0.25. What competing quote does it imply?
\end{exercise}

\begin{exercise}[$\star\star$]\label{exo:s2:flow-market-making-at-a-bank:4}
Why does no flat price match pricing by client?
\end{exercise}

\begin{exercise}[$\star\star$]\label{exo:s2:flow-market-making-at-a-bank:5}
Why does a small skew raise the \omterm{def:s2:flow-market-making-at-a-bank:franchise}{franchise P\&L} while a large one lowers it?
\end{exercise}

\begin{exercise}[$\star\star$]\label{exo:s2:flow-market-making-at-a-bank:6}
Why are internalisation costs lower for dealers whose clients are more price-sensitive?
\end{exercise}

\begin{exercise}[$\star\star\star$]\label{exo:s2:flow-market-making-at-a-bank:7}
\emph{Coding.} Run \texttt{policy} with a flat half-spread of 0.8 basis points. What happens to trades, marks and the \omterm{def:s2:flow-market-making-at-a-bank:franchise}{franchise P\&L}?
\end{exercise}

\begin{exercise}[$\star\star\star$]\label{exo:s2:flow-market-making-at-a-bank:8}
\emph{Find the flaw.} ``Our informed clients generate most of our volume growth; we should tighten their spreads to win more of it.''
\end{exercise}

\section{Problem: Pricing the Client}

\begin{problem}[{Weekend problem --- flow market making}]\label{pb:s2:flow-market-making-at-a-bank:1}
The chapter's synthetic desk and the public record.

\medskip\noindent\textbf{Part I --- The franchise.}
\begin{enumerate}
\item Define the \omterm{def:s2:flow-market-making-at-a-bank:franchise}{franchise P\&L} and its parts.
\item Describe the clients and why the mid path does not depend on the desk.
\item What do the mark-outs and hit ratios show after the flat history?
\item What did Oomen's price signatures show?
\end{enumerate}

\medskip\noindent\textbf{Part II --- Tiering.}
\begin{enumerate}[resume]
\item Define \omterm{def:s2:flow-market-making-at-a-bank:tiering}{client tiering}.
\item How is each client priced?
\item Give the policy table.
\item Give the capture and net per trade by type.
\end{enumerate}

\medskip\noindent\textbf{Part III --- Skews and internalisation.}
\begin{enumerate}[resume]
\item Define an \omterm{def:s2:flow-market-making-at-a-bank:skew}{axe-driven skew}.
\item What does the skew do to P\&L, inventory and internalisation?
\item What did Butz and Oomen find?
\item What is last look, and what did Oomen find?
\end{enumerate}

\medskip\noindent\textbf{Part IV --- The verdict.}
\begin{enumerate}[resume]
\item State the \emph{named result}: the desk's spread capture by tier and the share of risk internalised.
\item Which assumption flatters pricing by client most?
\item Why does a low flat price lose?
\item How would you estimate a client's win curve honestly?
\item What would you do with a client whose mark-out rises?
\item Which strategy file carries the most inventory risk?
\item How does this chapter relate to chapter~15's flow strategies?
\item In one sentence: what does a flow desk sell?
\end{enumerate}
\end{problem}

\section{Interview questions}

\begin{interviewq}[$\star$ \iqroles{trader}]\label{iq:s2:flow-market-making-at-a-bank:1}
What is a mark-out, and how would you use it to price a client?
\end{interviewq}

\begin{interviewq}[$\star\star$ \iqroles{researcher}]\label{iq:s2:flow-market-making-at-a-bank:2}
You see only the requests you won. How do you estimate a client's sensitivity to your price?
\end{interviewq}

\begin{interviewq}[$\star\star$ \iqroles{trader}]\label{iq:s2:flow-market-making-at-a-bank:3}
You are long a large position and a client asks for a two-way price. What do you quote?
\end{interviewq}

\begin{interviewq}[$\star\star$ \iqroles{risk}]\label{iq:s2:flow-market-making-at-a-bank:4}
How would you set the inventory limit of an internalising desk?
\end{interviewq}

\begin{interviewq}[$\star\star$ \iqroles{developer}]\label{iq:s2:flow-market-making-at-a-bank:5}
Design the service that computes client mark-outs in real time for a pricing engine.
\end{interviewq}

\begin{interviewq}[$\star\star\star$ \iqroles{researcher}]\label{iq:s2:flow-market-making-at-a-bank:6}
With a logistic win probability $p(h) = 1/(1 + e^{(h - c)/w})$ and a mark-out $m$, show that the profit-maximising half-spread satisfies
$h - m = w / (1 - p(h))$.
\end{interviewq}
